\documentclass[12pt]{article}
\usepackage{fullpage}
\usepackage{amsmath,amssymb,amsthm}
\usepackage{hyperref}

\newtheorem{theorem}{Theorem}
\newtheorem{lemma}{Lemma}
\newtheorem{proposition}{Proposition}
\newtheorem{corollary}{Corollary}
\theoremstyle{definition}
\newtheorem{definition}{Definition}
\newtheorem{remark}{Remark}

\DeclareMathOperator{\sgn}{sgn}

\begin{document}

\begin{center}
\textbf{\large A rationally generating family of $4k$--manifolds, by direct linear
algebra}\\[2pt]
\textit{(self-contained resolution of the ``generation gap'', issue rm-01541-2)}
\end{center}

\bigskip

\noindent
\textbf{Scope.}
This section supplies, by elementary means, the linear-algebra step that the standard
proof of the Hirzebruch signature theorem usually imports from cobordism theory
(Thom's computation $\Omega^{SO}_\ast\otimes\mathbb{Q}=
\mathbb{Q}[[\mathbb{CP}^2],[\mathbb{CP}^4],\dots]$).  We never use the cobordism ring,
the Pontryagin--Thom construction, or any spectral/index-theoretic input.  Everything
below is a finite computation with Pontryagin classes of explicit manifolds and a
determinant estimate.

\medskip

\noindent
\textbf{What is proved here.}
For each $k\ge 1$ we exhibit a finite, explicit family of closed oriented
$4k$--manifolds whose Pontryagin-number vectors are linearly independent and span the
full space of Pontryagin numbers in degree $4k$ (Theorem~\ref{thm:span}); we compute in
closed form the value on each family member of both functionals occurring in the
signature theorem, the signature $\tau$ and the $L$-number
$L[M]=\langle L(p_\bullet(M)),[M]\rangle$ (Proposition~\ref{prop:values}); and we deduce
that these two functionals coincide on degree-$4k$ Pontryagin numbers
(Corollary~\ref{cor:agree}).  The only external inputs are: (i) the multiplicativity of
both $\tau$ and $L[\cdot]$ under products, and (ii) the fact (established elsewhere in
the proof) that $\tau$ is expressible through Pontryagin numbers in each degree.  Both
(i) and (ii) are recalled precisely below; the genuinely combinatorial content -- the
spanning -- is proved here from scratch.

\bigskip


\section{Foundations: cohomology, intersection form, signature, and Pontryagin
classes of $\mathbb{CP}^{n}$}\label{sec:found}

The whole normalization of the signature theorem rests on two facts about complex
projective space: the structure of its cohomology ring (which pins down the intersection
form and hence the signature) and its total Pontryagin class.  We prove both from
scratch here, quoting only standard textbook results (the cohomology ring of
$\mathbb{CP}^n$, the Euler sequence, the splitting principle, and the Whitney product
formula) at the level at which they are routinely established.  All later sections refer
back to this one.

\subsection*{Cohomology ring and the fundamental class}

\begin{lemma}[Cohomology ring of $\mathbb{CP}^n$]\label{lem:cpring}
There is a class $h\in H^2(\mathbb{CP}^n;\mathbb{Z})$ (the first Chern class of the
hyperplane bundle $\mathcal{O}(1)$, equivalently the Poincar\'e dual of a hyperplane
$\mathbb{CP}^{n-1}\subset\mathbb{CP}^n$) such that
\[
H^*(\mathbb{CP}^n;\mathbb{Z})=\mathbb{Z}[h]/(h^{n+1}),
\qquad
H^*(\mathbb{CP}^n;\mathbb{R})=\mathbb{R}[h]/(h^{n+1}),
\]
so that $H^{2i}(\mathbb{CP}^n;\mathbb{R})=\mathbb{R}\cdot h^i$ is one-dimensional for
$0\le i\le n$ and $H^{\mathrm{odd}}=0$.  With the complex (hence canonical) orientation of
$\mathbb{CP}^n$, the fundamental class $[\mathbb{CP}^n]$ is normalized by
\[
\big\langle h^{\,n},[\mathbb{CP}^n]\big\rangle=1.
\tag{$\ast$}\label{eq:topdeg}
\]
\end{lemma}

\begin{proof}
The ring structure $H^*(\mathbb{CP}^n;\mathbb{Z})=\mathbb{Z}[h]/(h^{n+1})$ with
$h$ a generator of $H^2$ is the standard computation of the cohomology of complex
projective space (e.g.\ from the cell structure with one cell in each even dimension
$0,2,\dots,2n$, the cup-product structure being determined by the inclusion
$\mathbb{CP}^{n-1}\hookrightarrow\mathbb{CP}^n$ inducing an isomorphism in degrees
$\le 2n-2$); tensoring with $\mathbb{R}$ gives the real statement.  Thus each
$H^{2i}(\mathbb{CP}^n;\mathbb{R})$ is the line spanned by $h^i$, and the odd cohomology
vanishes.

For \eqref{eq:topdeg}: $h$ is by definition $c_1(\mathcal{O}(1))$, the Poincar\'e dual of
a hyperplane $H_0=\mathbb{CP}^{n-1}$.  Inductively, $h^i$ is Poincar\'e dual to a linear
subspace $\mathbb{CP}^{n-i}$ (the transverse intersection of $i$ generic hyperplanes), and
$h^n$ is Poincar\'e dual to a single point $\mathbb{CP}^0$, taken with positive sign
because distinct complex hyperplanes always meet positively (complex intersections carry
the canonical orientation).  Pairing the Poincar\'e dual of a positively oriented point
against the fundamental class gives $+1$, i.e.\ $\langle h^n,[\mathbb{CP}^n]\rangle=1$.
(Equivalently, one normalizes the orientation so that the generator $h^n\in H^{2n}$
evaluates to $+1$ on $[\mathbb{CP}^n]$; this is the standard complex orientation.)
\end{proof}

\begin{proposition}[Intersection form and signature of $\mathbb{CP}^{2k}$]\label{prop:sig}
For every $k\ge 0$ the middle cohomology of $\mathbb{CP}^{2k}$ is
\[
H^{2k}(\mathbb{CP}^{2k};\mathbb{R})=\mathbb{R}\cdot h^{k},
\]
one-dimensional, and the intersection form on it is
\[
(\alpha,\beta)\longmapsto\langle\alpha\smile\beta,[\mathbb{CP}^{2k}]\rangle,
\qquad
(h^k,h^k)\longmapsto\big\langle h^{k}\smile h^{k},[\mathbb{CP}^{2k}]\big\rangle
=\big\langle h^{2k},[\mathbb{CP}^{2k}]\big\rangle=1.
\]
Thus the intersection form is the rank-one positive-definite form $(1)$, and
\[
\tau(\mathbb{CP}^{2k})=1.
\]
More generally $\tau(\mathbb{CP}^m)=1$ for every even $m$ (apply the above with
$m=2k$), while for odd $m$ the middle degree $m$ is odd, $H^m(\mathbb{CP}^m;\mathbb{R})=0$,
and the signature is the empty signature.
\end{proposition}

\begin{proof}
By Lemma~\ref{lem:cpring} the dimension of $\mathbb{CP}^{2k}$ is $4k$, its middle real
cohomology $H^{2k}$ is the line $\mathbb{R}\cdot h^k$, and
$\langle h^{2k},[\mathbb{CP}^{2k}]\rangle=1$ by \eqref{eq:topdeg}.  Hence on the basis
vector $h^k$ the symmetric bilinear intersection form takes the value
$\langle h^k\smile h^k,[\mathbb{CP}^{2k}]\rangle=\langle h^{2k},[\mathbb{CP}^{2k}]\rangle
=1>0$.  A rank-one symmetric form with a positive diagonal entry is positive definite, so
its signature (number of positive eigenvalues minus number of negative ones) is
$1-0=1$.  Therefore $\tau(\mathbb{CP}^{2k})=1$.
\end{proof}

\subsection*{Total Pontryagin class}

We write $\underline{\mathbb{C}}$ for the trivial complex line bundle, $\mathcal{O}(1)$
for the hyperplane (dual tautological) bundle on $\mathbb{CP}^n$, and recall that for a
complex vector bundle $E$ the (real) Pontryagin classes are defined by
$p_i(E_{\mathbb{R}})=(-1)^i c_{2i}(E\otimes_{\mathbb{R}}\mathbb{C})$, and that for the
realification of a complex bundle one has the convenient identity, valid via the
splitting principle,
\[
p(E_{\mathbb{R}})=\prod_j (1+x_j^2),
\tag{$\dagger$}\label{eq:pfromc}
\]
where $x_1,\dots,x_r$ are the Chern roots of $E$ (so $c(E)=\prod_j(1+x_j)$).  Identity
\eqref{eq:pfromc} is the standard relation between Pontryagin and Chern classes of a
complex bundle; we record its short derivation below for completeness.

\begin{lemma}[Pontryagin from Chern via roots]\label{lem:pfromc}
For a complex vector bundle $E$ with Chern roots $x_j$, the total Pontryagin class of the
underlying real bundle is $p(E_{\mathbb{R}})=\prod_j(1+x_j^2)$.
\end{lemma}

\begin{proof}
By the splitting principle we may assume $E=\bigoplus_j \ell_j$ with $\ell_j$ complex line
bundles, $x_j=c_1(\ell_j)$.  The complexification of the realification satisfies
$\ell_{j,\mathbb{R}}\otimes_{\mathbb{R}}\mathbb{C}\cong \ell_j\oplus\bar\ell_j$, with Chern
roots $x_j$ and $-x_j$ (since $c_1(\bar\ell_j)=-c_1(\ell_j)$).  Hence
$c(E_{\mathbb{R}}\otimes\mathbb{C})=\prod_j(1+x_j)(1-x_j)=\prod_j(1-x_j^2)$, and
\[
p(E_{\mathbb{R}})=\sum_i (-1)^i c_{2i}(E_{\mathbb{R}}\otimes\mathbb{C})
=\prod_j\big(1+x_j^2\big),
\]
the last equality because replacing $x_j^2\mapsto -x_j^2$ (equivalently inserting the sign
$(-1)^i$ on the degree-$2i$ part) turns $\prod_j(1-x_j^2)$ into $\prod_j(1+x_j^2)$.
\end{proof}

\begin{proposition}[Total Pontryagin class of $\mathbb{CP}^n$]\label{prop:cppont}
With $h=c_1(\mathcal{O}(1))$ as above,
\[
c(T\mathbb{CP}^n)=(1+h)^{n+1},
\qquad
p(T\mathbb{CP}^n)=(1+h^2)^{\,n+1},
\qquad
p_i(\mathbb{CP}^n)=\binom{n+1}{i}h^{2i}.
\]
\end{proposition}

\begin{proof}
\emph{Step 1: the Euler sequence and $T\mathbb{CP}^n\oplus\underline{\mathbb{C}}\cong
(n+1)\,\mathcal{O}(1)$.}
Let $\gamma=\mathcal{O}(-1)$ be the tautological line bundle, whose fiber over a line
$\ell\subset\mathbb{C}^{n+1}$ is $\ell$ itself, and let $\underline{\mathbb{C}^{n+1}}$ be
the trivial bundle of rank $n+1$.  The Euler sequence is the short exact sequence of
holomorphic bundles
\[
0\longrightarrow \underline{\mathbb{C}}\longrightarrow
\operatorname{Hom}(\gamma,\underline{\mathbb{C}^{n+1}})
\longrightarrow T\mathbb{CP}^n\longrightarrow 0,
\]
where the first map sends $1$ to the inclusion $\gamma\hookrightarrow
\underline{\mathbb{C}^{n+1}}$, and the surjection identifies
$\operatorname{Hom}(\gamma,\underline{\mathbb{C}^{n+1}}/\gamma)
\cong T\mathbb{CP}^n$ (the standard description of the holomorphic tangent space at $\ell$
as $\operatorname{Hom}(\ell,\mathbb{C}^{n+1}/\ell)$).  Since
$\operatorname{Hom}(\gamma,\underline{\mathbb{C}^{n+1}})\cong
\gamma^\vee\otimes\underline{\mathbb{C}^{n+1}}\cong (n+1)\,\gamma^\vee
=(n+1)\,\mathcal{O}(1)$, and every short exact sequence of complex vector bundles over a
paracompact base splits (smoothly), we obtain
\[
T\mathbb{CP}^n\oplus\underline{\mathbb{C}}\cong(n+1)\,\mathcal{O}(1).
\]

\emph{Step 2: Chern class.}  Chern classes are stable (a trivial summand does not change
them) and multiplicative (Whitney formula), and $c(\mathcal{O}(1))=1+h$ by definition of
$h$.  Therefore
\[
c(T\mathbb{CP}^n)=c\big(T\mathbb{CP}^n\oplus\underline{\mathbb{C}}\big)
=c\big((n+1)\mathcal{O}(1)\big)=\big(c(\mathcal{O}(1))\big)^{n+1}=(1+h)^{n+1}.
\]
In particular all $n+1$ Chern roots of $T\mathbb{CP}^n$ equal $h$.

\emph{Step 3: Pontryagin class.}  By Lemma~\ref{lem:pfromc}, replacing each Chern root
$x_j=h$ by $1+x_j^2=1+h^2$ gives
\[
p(T\mathbb{CP}^n)=\prod_{j=1}^{n+1}(1+h^2)=(1+h^2)^{n+1}
=\sum_{i\ge0}\binom{n+1}{i}h^{2i},
\]
so $p_i(\mathbb{CP}^n)=\binom{n+1}{i}h^{2i}$ (a class in $H^{4i}$; it vanishes once
$2i>n$, consistent with $h^{n+1}=0$).  Pontryagin classes are likewise unchanged by the
trivial summand $\underline{\mathbb{C}}$, so this is indeed $p(T\mathbb{CP}^n)$.
\end{proof}

These two propositions are exactly the inputs used below: Proposition~\ref{prop:sig}
supplies $\tau(\mathbb{CP}^{2k})=1$ and Proposition~\ref{prop:cppont} supplies the formula
$p(T\mathbb{CP}^n)=(1+h^2)^{n+1}$, on which the Pontryagin-number computations and the
$L$-number evaluation rest.

\bigskip

\section{The space of Pontryagin numbers and the two functionals}

Throughout, $k\ge 1$ is fixed and $P(k)$ denotes the set of partitions of $k$; we write
$p(k)=\#P(k)$.  A partition is written $\lambda=(\lambda_1\ge\lambda_2\ge\cdots\ge
\lambda_r\ge 1)$ with $|\lambda|=\sum_i\lambda_i=k$ and length $\ell(\lambda)=r$.

\begin{definition}[Pontryagin numbers]\label{def:pn}
Let $M$ be a closed oriented smooth $4k$--manifold with Pontryagin classes
$p_i=p_i(TM)\in H^{4i}(M;\mathbb{Q})$.  For $\lambda\in P(k)$ put
\[
p_\lambda[M]\;:=\;\big\langle\, p_{\lambda_1}\smile p_{\lambda_2}\smile\cdots\smile
p_{\lambda_r}\,,\,[M]\,\big\rangle\;\in\;\mathbb{Q}.
\]
The \emph{Pontryagin vector} of $M$ is
$\mathbf{p}[M]=\big(p_\lambda[M]\big)_{\lambda\in P(k)}\in\mathbb{Q}^{P(k)}$.
\end{definition}

The $p(k)$ monomials $p_\lambda$ form a $\mathbb{Q}$-basis of the degree-$4k$ part of
the free graded-commutative polynomial algebra $\mathbb{Q}[p_1,p_2,\dots]$ on generators
$p_i$ of degree $4i$.  Consequently the degree-$4k$ component of any multiplicative
sequence (in particular Hirzebruch's $L$-polynomial) is a $\mathbb{Q}$-linear
combination of the $p_\lambda$; the $L$-number is therefore a fixed linear functional of
the Pontryagin vector:
\[
L[M]\;=\;\big\langle L(p_\bullet(M)),[M]\big\rangle\;=\;
\sum_{\lambda\in P(k)} c_\lambda\, p_\lambda[M],
\qquad c_\lambda\in\mathbb{Q}\ \text{fixed},
\tag{1}\label{eq:Llinear}
\]
where the constants $c_\lambda$ are the coefficients of the degree-$4k$ part of the
$L$-series (the precise values of the $c_\lambda$ are irrelevant for the present step).

The signature theorem is the assertion that the integer-valued functional $M\mapsto
\tau(M)$ equals the functional \eqref{eq:Llinear}.  In the elementary proof the
following two facts are established by independent (non-cobordism) arguments, and we take
them as the interface to the rest of the proof:

\begin{itemize}
\item[\textbf{(I)}] \emph{Multiplicativity.} For closed oriented manifolds $M,N$ of
dimensions divisible by $4$,
\[
\tau(M\times N)=\tau(M)\,\tau(N),\qquad L[M\times N]=L[M]\cdot L[N].
\tag{2}\label{eq:mult}
\]
(For $L[\cdot]$ this is the multiplicativity of multiplicative sequences together with
the Whitney formula $p(T(M\times N))=p(TM)\times p(TN)$ and
$[M\times N]=[M]\times[N]$; for $\tau$ it is the elementary fact that the intersection
form of a product is the tensor product of the intersection forms, whose signature is
the product of the signatures.)
\item[\textbf{(II)}] \emph{$\tau$ is a Pontryagin functional.} In each degree $4k$ there
exist (a priori unknown) constants $a_\lambda\in\mathbb{Q}$ with
$\tau(M)=\sum_{\lambda\in P(k)}a_\lambda\, p_\lambda[M]$ for every closed oriented
$4k$--manifold $M$.
\end{itemize}

\noindent
Fact \eqref{eq:mult} is elementary and recalled above.  Fact (II) is the
non-combinatorial heart of the theorem and is supplied elsewhere in the proof; it does
\emph{not} rely on the present section.  The role of the present section is to show that
(I), (II), and the value computations below force $a_\lambda=c_\lambda$ for all
$\lambda$, i.e.\ $\tau=L[\cdot]$.

\bigskip

\section{The generating family and its Pontryagin numbers}

\begin{definition}[The family]\label{def:family}
For $\mu=(\mu_1\ge\cdots\ge\mu_r)\in P(k)$ set
\[
M_\mu\;:=\;\mathbb{CP}^{2\mu_1}\times\mathbb{CP}^{2\mu_2}\times\cdots\times
\mathbb{CP}^{2\mu_r},
\]
a closed oriented manifold of real dimension $\sum_s 4\mu_s=4k$.  This is a family of
exactly $p(k)$ manifolds, indexed by $P(k)$.
\end{definition}

We now compute all Pontryagin numbers of $M_\mu$ in closed form.  Recall the standard
fact, proved purely from the Euler sequence and the splitting principle: with
$h\in H^2(\mathbb{CP}^n;\mathbb{Z})$ the hyperplane class,
\[
p(T\mathbb{CP}^n)\;=\;(1+h^2)^{\,n+1},
\qquad\text{so}\qquad
p_i(\mathbb{CP}^n)=\binom{n+1}{i}h^{2i},
\tag{3}\label{eq:cppont}
\]
and $\langle h^n,[\mathbb{CP}^n]\rangle=1$.  Equation \eqref{eq:cppont} is precisely
Proposition~\ref{prop:cppont}, and $\langle h^n,[\mathbb{CP}^n]\rangle=1$ is
\eqref{eq:topdeg} of Lemma~\ref{lem:cpring}; both are proved in
Section~\ref{sec:found}.

Introduce, for each factor $s$ of $M_\mu$, the degree-$4$ class
\[
t_s\;:=\;h_s^2\in H^4(M_\mu;\mathbb{Q}),
\qquad s=1,\dots,r,
\]
where $h_s$ is the hyperplane class pulled back from the $s$-th factor
$\mathbb{CP}^{2\mu_s}$.  By \eqref{eq:cppont} and the Whitney formula, the total
Pontryagin class of $M_\mu$ is
\[
p(TM_\mu)\;=\;\prod_{s=1}^{r}(1+t_s)^{\,2\mu_s+1},
\tag{4}\label{eq:prodp}
\]
and the fundamental cohomology class is dual to $\prod_s h_s^{2\mu_s}=\prod_s t_s^{\mu_s}$,
so for any polynomial $Q$ in the $t_s$,
\[
\langle Q,[M_\mu]\rangle\;=\;\big[\textstyle\prod_s t_s^{\mu_s}\big]\,Q,
\tag{5}\label{eq:pair}
\]
the coefficient of the monomial $\prod_s t_s^{\mu_s}$ in $Q$.

It is convenient to work with the \emph{power-sum Pontryagin numbers}.  Let
\[
\mathfrak{p}_n\;:=\;\sum_{j} y_j^{\,n}
\]
denote the $n$-th power sum of the formal Pontryagin roots $y_j$ (so that, as usual,
$1+p_1+p_2+\cdots=\prod_j(1+y_j)$ and the $\mathfrak{p}_n$ are the Newton polynomials in
the $p_i$ with \emph{rational} coefficients).  For a partition
$\lambda=(\lambda_1,\dots,\lambda_r)\in P(k)$ define
\[
\mathfrak{p}_\lambda\;:=\;\mathfrak{p}_{\lambda_1}\mathfrak{p}_{\lambda_2}\cdots
\mathfrak{p}_{\lambda_r},
\qquad
\mathfrak{p}_\lambda[M]\;:=\;\langle\mathfrak{p}_\lambda,[M]\rangle.
\]
Because the Newton relations express the $\mathfrak{p}_n$ in terms of the $p_i$ and vice
versa by an \emph{invertible} (unitriangular up to nonzero scalars) change of basis over
$\mathbb{Q}$, the families $\{\mathfrak{p}_\lambda\}_{\lambda\in P(k)}$ and
$\{p_\lambda\}_{\lambda\in P(k)}$ span the same degree-$4k$ subspace and
\[
\operatorname{span}_\mathbb{Q}\{\mathbf{p}[M_\mu]:\mu\in P(k)\}
=\mathbb{Q}^{P(k)}
\iff
\det\big(\mathfrak{p}_\lambda[M_\mu]\big)_{\lambda,\mu\in P(k)}\neq 0.
\tag{6}\label{eq:reduce}
\]

For $M_\mu$, the formal Pontryagin roots are, by \eqref{eq:prodp}, the classes $t_s$ each
occurring with multiplicity $2\mu_s+1$.  Hence
\[
\mathfrak{p}_n(M_\mu)=\sum_{s=1}^{r}(2\mu_s+1)\,t_s^{\,n},
\]
and therefore, using \eqref{eq:pair},
\[
\mathfrak{p}_\lambda[M_\mu]
=\Big[\textstyle\prod_{s} t_s^{\mu_s}\Big]\
\prod_{i=1}^{\ell(\lambda)}\Big(\sum_{s=1}^{\ell(\mu)}(2\mu_s+1)\,t_s^{\lambda_i}\Big).
\tag{7}\label{eq:closed}
\]
Formula \eqref{eq:closed} is an explicit, finite, integer-valued closed form for every
entry of the matrix in \eqref{eq:reduce}.

\bigskip

\section{Triangularity and nonsingularity}

We now prove that the matrix
$A:=\big(\mathfrak{p}_\lambda[M_\mu]\big)_{\lambda,\mu\in P(k)}$ is nonsingular, by a
direct combinatorial argument (no cobordism, no symmetric-function black boxes beyond the
elementary Newton change of basis already invoked).

Recall the \emph{dominance order} on $P(k)$: for $\alpha,\beta\in P(k)$ written in
weakly decreasing order and padded with zeros to equal length,
\[
\alpha\preceq\beta\iff
\sum_{i\le j}\alpha_i\le\sum_{i\le j}\beta_i\quad\text{for all }j.
\]

\begin{lemma}[Combinatorial expansion of the entries]\label{lem:expand}
For $\lambda,\mu\in P(k)$,
\[
\mathfrak{p}_\lambda[M_\mu]
=\sum_{f}\ \prod_{i=1}^{\ell(\lambda)}\big(2\mu_{f(i)}+1\big),
\tag{8}\label{eq:expand}
\]
where the sum runs over all functions $f:\{1,\dots,\ell(\lambda)\}\to\{1,\dots,\ell(\mu)\}$
such that
\[
\sum_{i:\,f(i)=s}\lambda_i=\mu_s\quad\text{for every }s=1,\dots,\ell(\mu).
\tag{9}\label{eq:balance}
\]
In particular every entry $\mathfrak{p}_\lambda[M_\mu]$ is a nonnegative integer, and it
is nonzero if and only if at least one such function $f$ exists.
\end{lemma}

\begin{proof}
Expand the product in \eqref{eq:closed}.  Each factor indexed by $i$ contributes one term
$(2\mu_{s}+1)\,t_{s}^{\lambda_i}$, i.e.\ a choice $f(i)=s\in\{1,\dots,\ell(\mu)\}$ of which
summand is taken.  A choice of $f$ produces the monomial
$\prod_{i}(2\mu_{f(i)}+1)\,t_{f(i)}^{\lambda_i}
=\big(\prod_i(2\mu_{f(i)}+1)\big)\,\prod_{s}t_s^{\sum_{i:f(i)=s}\lambda_i}$.
By \eqref{eq:pair} the coefficient of $\prod_s t_s^{\mu_s}$ is the sum of
$\prod_i(2\mu_{f(i)}+1)$ over precisely those $f$ for which the exponent of each $t_s$
equals $\mu_s$, which is condition \eqref{eq:balance}.  All weights
$2\mu_s+1\ge 3>0$, so the sum is a nonnegative integer, zero exactly when no $f$ satisfies
\eqref{eq:balance}.
\end{proof}

\begin{lemma}[Coarsening dominates]\label{lem:merge}
Let $\lambda\in P(k)$ and let $\{B_1,\dots,B_m\}$ be a partition of the index set
$\{1,\dots,\ell(\lambda)\}$ into nonempty blocks.  Define $\mu_s:=\sum_{i\in B_s}\lambda_i$
for $s=1,\dots,m$; then $\mu=(\mu_1,\dots,\mu_m)$ is a partition of $k$ (after weakly
decreasing rearrangement) and $\lambda\preceq\mu$ in the dominance order.
\end{lemma}

\begin{proof}
First, a single \emph{pairwise merge}.  Let $\nu$ be a partition and let $a\ge b\ge 1$ be
two of its parts; let $\nu'$ be obtained by deleting one occurrence each of $a$ and $b$ and
inserting the single part $a+b$ (so $|\nu'|=|\nu|$).  We claim $\nu\preceq\nu'$.  Write both
partitions in weakly decreasing order.  Passing from $\nu$ to $\nu'$ replaces two parts
$a,b$ by the larger part $a+b$ and one fewer total part.  Consider the partial sums
$S_j(\cdot):=\sum_{i\le j}(\cdot)^{\downarrow}_i$.  For any threshold value, the multiset of
parts of $\nu'$ that are $\ge$ that threshold is obtained from that of $\nu$ by possibly
removing $a$ and $b$ and possibly adding $a+b\ge a\ge b$.  Concretely: the part $a+b$ in
$\nu'$ is at least as large as either of $a,b$, so for every $j$ the $j$ largest parts of
$\nu'$ have sum at least the $j$ largest parts of $\nu$.  Indeed, fix $j$.  Let $T$ be the
multiset of the $j$ largest parts of $\nu$ (sum $S_j(\nu)$).  Build a sub-multiset $T'$ of
the parts of $\nu'$ with $|T'|\le j$ and $\sum T'\ge \sum T$ as follows: keep every element
of $T$ that is still a part of $\nu'$ (i.e.\ all elements of $T$ except possibly the deleted
copies of $a$ and $b$); if $T$ contained at least one of $a,b$, replace those (one or two)
deleted elements by the single part $a+b$, which is $\ge$ their sum if both were present and
$\ge$ the single one otherwise.  This $T'$ has $|T'|\le j$ and $\sum T'\ge\sum T=S_j(\nu)$.
Since the $j$ largest parts of $\nu'$ maximize the sum over all sub-multisets of size
$\le j$, we get $S_j(\nu')\ge\sum T'\ge S_j(\nu)$.  As $j$ was arbitrary and total sums
agree, $\nu\preceq\nu'$.

Now the general statement.  The passage from $\lambda$ to $\mu$ replaces each block
$B_s=\{i_1,\dots,i_{|B_s|}\}$ by the single part $\sum_{i\in B_s}\lambda_i$, which is the
result of $|B_s|-1$ successive pairwise merges within that block; performing all such merges
across all blocks expresses $\mu$ as the end of a finite chain
$\lambda=\nu^{(0)}\preceq\nu^{(1)}\preceq\cdots\preceq\nu^{(N)}=\mu$, each step a pairwise
merge.  By the previous paragraph each step is a dominance increase, and $\preceq$ is
transitive, so $\lambda\preceq\mu$.
\end{proof}

\begin{lemma}[Triangularity]\label{lem:tri}
If $\mathfrak{p}_\lambda[M_\mu]\neq 0$ then $\lambda\preceq\mu$.  Moreover
$\mathfrak{p}_\mu[M_\mu]>0$ for every $\mu\in P(k)$.
\end{lemma}

\begin{proof}
Suppose $\mathfrak{p}_\lambda[M_\mu]\neq 0$.  By Lemma~\ref{lem:expand} there is a
function $f:\{1,\dots,\ell(\lambda)\}\to\{1,\dots,\ell(\mu)\}$ satisfying
\eqref{eq:balance}.  The fibers $B_s:=f^{-1}(s)$ partition $\{1,\dots,\ell(\lambda)\}$; by
\eqref{eq:balance} each is nonempty (if some $B_s$ were empty then $\mu_s=0$, contradicting
$\mu_s\ge 1$), and $\mu_s=\sum_{i\in B_s}\lambda_i$.  Thus $\mu$ is the coarsening of
$\lambda$ associated to the blocks $\{B_s\}$, so by Lemma~\ref{lem:merge},
$\lambda\preceq\mu$.

For the diagonal, take $\lambda=\mu$ and let $f=\mathrm{id}$ (after fixing any ordering of
the parts of $\mu$); then $\sum_{i:f(i)=s}\mu_i=\mu_s$, so $f$ satisfies \eqref{eq:balance}.
By \eqref{eq:expand} the entry $\mathfrak{p}_\mu[M_\mu]$ is a sum of strictly positive terms
(one for each such $f$), hence $\mathfrak{p}_\mu[M_\mu]>0$.
\end{proof}

\begin{theorem}[Spanning]\label{thm:span}
For every $k\ge 1$ the $p(k)$ manifolds $\{M_\mu:\mu\in P(k)\}$ have linearly independent
Pontryagin vectors, and these vectors span $\mathbb{Q}^{P(k)}$.  Equivalently, the only
$\mathbb{Q}$-linear functional of degree-$4k$ Pontryagin numbers that vanishes on every
$M_\mu$ is the zero functional.
\end{theorem}

\begin{proof}
Order $P(k)$ by any linear extension of the dominance order $\preceq$ (such an extension
exists since $\preceq$ is a partial order on the finite set $P(k)$).  By
Lemma~\ref{lem:tri}, in this ordering the matrix
$A=\big(\mathfrak{p}_\lambda[M_\mu]\big)_{\lambda,\mu}$ is triangular: its $(\lambda,\mu)$
entry vanishes whenever $\lambda\not\preceq\mu$, i.e.\ whenever $\lambda$ comes strictly
after $\mu$ in the chosen ordering.  Its diagonal entries
$\mathfrak{p}_\mu[M_\mu]>0$ are nonzero.  Hence
\[
\det A=\prod_{\mu\in P(k)}\mathfrak{p}_\mu[M_\mu]\;>\;0,
\]
so $A$ is nonsingular.  By the equivalence \eqref{eq:reduce}, the Pontryagin vectors
$\{\mathbf{p}[M_\mu]\}$ are linearly independent and span $\mathbb{Q}^{P(k)}$.  The final
sentence is the dual statement: a linear functional vanishing on a spanning set is zero.
\end{proof}

\begin{remark}
The determinant is in fact a product of positive integers; for small $k$ one finds
$\det A=3,\,90,\,17010,\,1653372000,\dots$ for $k=1,2,3,4$.  The argument uses only the
explicit Pontryagin classes \eqref{eq:cppont}, the Whitney formula, the elementary Newton
change of basis between $p_\lambda$ and $\mathfrak{p}_\lambda$, and the combinatorial
triangularity of Lemma~\ref{lem:tri}.  At no point is $\Omega^{SO}_\ast\otimes\mathbb{Q}$,
the Pontryagin--Thom construction, or any generation statement from cobordism theory
invoked; the spanning is a self-contained determinant computation.
\end{remark}

\bigskip

\section{Closed-form values of the two functionals}

\begin{proposition}[Values on the family]\label{prop:values}
For every $k\ge 1$ and every $\mu\in P(k)$,
\[
\tau(M_\mu)=1
\qquad\text{and}\qquad
L[M_\mu]=1.
\]
\end{proposition}

\begin{proof}
\emph{Signature.}  By Proposition~\ref{prop:sig}, $\tau(\mathbb{CP}^{2k})=1$ for every
$k\ge0$; in particular $\tau(\mathbb{CP}^{2\mu_s})=1$ for each factor $s$.  By the
multiplicativity \eqref{eq:mult}, $\tau(M_\mu)=\prod_s\tau(\mathbb{CP}^{2\mu_s})=1$.

\emph{$L$-number.}  By multiplicativity \eqref{eq:mult} of the $L$-genus it suffices to
show $L[\mathbb{CP}^{2m}]=1$ for all $m\ge 1$.  By \eqref{eq:cppont} the formal
Pontryagin roots of $\mathbb{CP}^{2m}$ are $h^2$ with multiplicity $2m+1$; hence the
$L$-polynomial evaluated on $\mathbb{CP}^{2m}$ is the degree-$4m$ part of
\[
L(p_\bullet)\Big|_{\mathbb{CP}^{2m}}
=\Big(\frac{\sqrt{h^2}}{\tanh\sqrt{h^2}}\Big)^{2m+1}
=\Big(\frac{h}{\tanh h}\Big)^{2m+1},
\]
where $h/\tanh h=\sum_{j\ge 0}c_j h^{2j}$ is the even power series defining the $L$-genus.
Pairing with $[\mathbb{CP}^{2m}]$ extracts the coefficient of $h^{2m}$:
\[
L[\mathbb{CP}^{2m}]=\big[h^{2m}\big]\Big(\frac{h}{\tanh h}\Big)^{2m+1}
=\frac{1}{2\pi i}\oint \Big(\frac{h}{\tanh h}\Big)^{2m+1}\frac{dh}{h^{2m+1}}
=\frac{1}{2\pi i}\oint\frac{dh}{(\tanh h)^{2m+1}},
\]
the contour being a small circle about $h=0$.  Substitute $w=\tanh h$, so
$dw=(1-w^2)\,dh=\operatorname{sech}^2 h\,dh$ and $h=0\mapsto w=0$ with $dw/dh=1$ at the
origin; the substitution is a local biholomorphism near $0$, and
\[
\frac{1}{2\pi i}\oint\frac{dh}{(\tanh h)^{2m+1}}
=\frac{1}{2\pi i}\oint\frac{1}{w^{2m+1}}\,\frac{dw}{1-w^2}
=\big[w^{2m}\big]\frac{1}{1-w^2}
=\big[w^{2m}\big]\sum_{j\ge0}w^{2j}=1.
\]
Hence $L[\mathbb{CP}^{2m}]=1$, and by multiplicativity
$L[M_\mu]=\prod_s L[\mathbb{CP}^{2\mu_s}]=1$.
\end{proof}

\begin{corollary}[Agreement on Pontryagin numbers]\label{cor:agree}
For every $k\ge 1$, the two $\mathbb{Q}$-linear functionals of degree-$4k$ Pontryagin
numbers,
\[
M\longmapsto\tau(M)
\qquad\text{and}\qquad
M\longmapsto L[M]=\sum_\lambda c_\lambda\,p_\lambda[M],
\]
coincide; that is, $a_\lambda=c_\lambda$ for all $\lambda\in P(k)$, where $a_\lambda$ are
the constants of fact \textnormal{(II)}.  Consequently
\[
\tau(M)=L[M]=\big\langle L(p_\bullet(M)),[M]\big\rangle
\]
for every closed oriented smooth $4k$--manifold $M$.
\end{corollary}

\begin{proof}
By fact (II), $\tau(M)=\sum_\lambda a_\lambda p_\lambda[M]$ for fixed $a_\lambda$, while by
\eqref{eq:Llinear} $L[M]=\sum_\lambda c_\lambda p_\lambda[M]$.  Their difference
\[
\Delta(M):=\tau(M)-L[M]=\sum_{\lambda\in P(k)}(a_\lambda-c_\lambda)\,p_\lambda[M]
\]
is a $\mathbb{Q}$-linear functional of the degree-$4k$ Pontryagin numbers.  By
Proposition~\ref{prop:values}, $\Delta(M_\mu)=\tau(M_\mu)-L[M_\mu]=1-1=0$ for every
$\mu\in P(k)$.  Thus $\Delta$ vanishes on the spanning family of
Theorem~\ref{thm:span}, hence $\Delta\equiv 0$ as a functional, i.e.\
$a_\lambda=c_\lambda$ for all $\lambda$.  Therefore $\tau(M)=L[M]$ for every closed
oriented $4k$--manifold $M$.
\end{proof}

\bigskip

\begin{remark}[On the interface facts]
The corollary closes the signature theorem \emph{given} facts (I) and (II).  Fact (I) is
elementary and proved above.  Fact (II) -- that $\tau$ is, in each degree, a fixed linear
combination of Pontryagin numbers -- is the analytic/topological core of the theorem and
is established in the companion sections of this proof by elementary (non-cobordism)
means; it is logically independent of the present section.  The content delivered here,
and the precise object the cobordism-free argument was missing, is the
\emph{spanning family with closed-form values}: an explicit set of $p(k)$ manifolds whose
Pontryagin vectors are provably independent by a direct triangular-determinant
computation, on which both functionals are computed to equal $1$.  This is exactly what
the standard proof imports from $\Omega^{SO}_\ast\otimes\mathbb{Q}$, now supplied
elementarily.
\end{remark}

\bigskip

\section{The cobordism invariance of the signature, and the status of Fact
(II)}\label{sec:factII}

The corollary of the previous section closes the signature theorem \emph{given} the two
interface facts (I) and (II).  Fact (I) was proved above.  This section addresses Fact
(II):
\[
\textbf{(II)}\qquad
\tau(M)=\sum_{\lambda\in P(k)}a_\lambda\,p_\lambda[M]
\quad\text{for fixed rationals }a_\lambda\ \text{independent of }M .
\]
We first prove, completely and by elementary means (only Poincar\'e--Lefschetz duality
and linear algebra of the intersection form), the central structural fact on which Fact
(II) rests: \emph{the signature is an oriented-cobordism invariant}.  We then explain
precisely how Fact (II) follows from this together with one further input, and we state
that input correctly and identify exactly what it provides.

\subsection*{Signature and Pontryagin numbers vanish on boundaries}

By a \emph{compact oriented cobordism} we mean a compact smooth oriented
$(4k{+}1)$--manifold $W$ with $\partial W=M$, the boundary orientation being the induced
one.  No cobordism ring, classifying space, or Pontryagin--Thom construction is used; only
the manifold $W$ and Poincar\'e--Lefschetz duality enter.

\begin{theorem}[Thom; cobordism invariance of the signature]\label{thm:cobinv}
Let $M$ be a closed oriented smooth $4k$--manifold.  If $M=\partial W$ for some compact
oriented smooth $(4k{+}1)$--manifold $W$, then $\tau(M)=0$.
\end{theorem}

\begin{proof}
All cohomology is with real coefficients.  Write $n=4k$, so $\dim W=n+1=4k+1$, and let
$j^\ast\colon H^{2k}(W)\to H^{2k}(M)$ be the restriction induced by the inclusion
$j\colon M=\partial W\hookrightarrow W$.  Set
\[
A:=\operatorname{im}\big(j^\ast\colon H^{2k}(W)\to H^{2k}(M)\big)\subseteq H^{2k}(M).
\]
Let $b:=\dim_{\mathbb R} H^{2k}(M)$ and $r:=\dim_{\mathbb R}A$.  Denote by
$\langle\,\cdot,\cdot\,\rangle_M$ the intersection (cup-product) form on $H^{2k}(M)$,
which is nondegenerate by Poincar\'e duality on $M$.  We prove two claims:
\smallskip

\noindent\emph{Claim 1: $A$ is isotropic, i.e.\ $\langle x,y\rangle_M=0$ for all
$x,y\in A$.}\;
Let $x=j^\ast\xi$, $y=j^\ast\eta$ with $\xi,\eta\in H^{2k}(W)$.  By naturality of the cup
product under $j^\ast$, $x\smile y=j^\ast(\xi\smile\eta)$ in $H^{4k}(M)$, with
$\xi\smile\eta\in H^{4k}(W)$.  Let $[W,M]\in H_{4k+1}(W,M)$ be the relative fundamental
class, so that $\partial[W,M]=[M]\in H_{4k}(M)$.  Then $j_\ast[M]=j_\ast\partial[W,M]=0$
in $H_{4k}(W)$, since the composite $H_{4k+1}(W,M)\xrightarrow{\partial}H_{4k}(M)
\xrightarrow{j_\ast}H_{4k}(W)$ is zero by exactness of the long exact homology sequence of
the pair $(W,M)$.  Therefore, by the adjunction between $j^\ast$ and $j_\ast$ for the
Kronecker pairing,
\[
\big\langle x\smile y,[M]\big\rangle
=\big\langle j^\ast(\xi\smile\eta),[M]\big\rangle
=\big\langle \xi\smile\eta,\,j_\ast[M]\big\rangle
=\big\langle \xi\smile\eta,\,0\big\rangle=0 ,
\]
proving Claim 1.
\smallskip

\noindent\emph{Claim 2: $r=\dim A=\tfrac12\,b$.}\;
Consider the segment of the long exact sequence of the pair $(W,M)$ in degree $2k$,
together with Lefschetz duality.  By Poincar\'e--Lefschetz duality for the compact
oriented manifold $W$ with boundary, there is a nondegenerate pairing
$H^{i}(W)\times H^{n+1-i}(W,M)\to\mathbb R$ inducing isomorphisms
\[
H^{i}(W)\;\cong\;H_{\,n+1-i}(W,M)\;\cong\;H^{\,n+1-i}(W,M)^\ast .
\]
Specialise the long exact sequence of $(W,M)$,
\[
\cdots\to H^{2k}(W,M)\xrightarrow{\;p\;}H^{2k}(W)\xrightarrow{\;j^\ast\;}
H^{2k}(M)\xrightarrow{\;\delta\;}H^{2k+1}(W,M)\xrightarrow{\;p'\;}H^{2k+1}(W)\to\cdots .
\]
Lefschetz duality identifies $H^{2k}(W,M)\cong H^{2k+1}(W)^\ast$ and
$H^{2k}(W)\cong H^{2k+1}(W,M)^\ast$ (here $n+1-2k=2k+1$), and under these identifications
the map $p\colon H^{2k}(W,M)\to H^{2k}(W)$ is the transpose of
$p'\colon H^{2k+1}(W,M)\to H^{2k+1}(W)$.  Therefore
$\operatorname{rank}p=\operatorname{rank}p'$.  Exactness gives, on the one hand,
$r=\dim\operatorname{im}j^\ast=\dim H^{2k}(W)-\operatorname{rank}p$, and on the other hand
$\dim\operatorname{im}\delta=\dim H^{2k}(M)-r=b-r$ while
$\dim\operatorname{im}\delta=\dim\ker p'=\dim H^{2k+1}(W,M)-\operatorname{rank}p'$.  Now
Lefschetz duality also gives $\dim H^{2k}(W)=\dim H^{2k+1}(W,M)$ (dual spaces).  Combining,
\[
b-r=\dim\operatorname{im}\delta=\dim H^{2k+1}(W,M)-\operatorname{rank}p'
=\dim H^{2k}(W)-\operatorname{rank}p=r .
\]
Hence $b-r=r$, i.e.\ $r=\tfrac12 b$.  (In particular $b$ is even.)
\smallskip

\noindent\emph{Conclusion.}  The intersection form $\langle\,\cdot,\cdot\,\rangle_M$ is a
nondegenerate symmetric form on the $b$--dimensional space $H^{2k}(M)$, and by Claims 1--2
it possesses an isotropic subspace $A$ of dimension exactly $\tfrac12 b$.  A
nondegenerate symmetric real form of rank $b$ admits a totally isotropic subspace of
dimension $\tfrac12 b$ if and only if it has signature $0$: indeed, in an orthogonal
eigenbasis the form is $\operatorname{diag}(+1,\dots,+1,-1,\dots,-1)$ with $p$ plus signs
and $q=b-p$ minus signs, and the maximal dimension of an isotropic subspace equals
$\min(p,q)+0=\min(p,q)$ \big(Witt index $=\min(p,q)$\big); having an isotropic subspace of
dimension $\tfrac12 b=\tfrac{p+q}2$ forces $\min(p,q)\ge\tfrac{p+q}2$, hence $p=q$, i.e.\
$\tau(M)=p-q=0$.  Therefore $\tau(M)=0$.
\end{proof}

\begin{proposition}[Pontryagin numbers vanish on boundaries]\label{prop:pnbdy}
Let $M=\partial W$ as above.  Then every Pontryagin number of $M$ vanishes:
$p_\lambda[M]=0$ for all $\lambda\in P(k)$.
\end{proposition}

\begin{proof}
The tangent bundle satisfies $TW|_M\cong TM\oplus\nu$, where $\nu$ is the (trivial, rank
one) normal bundle of the boundary; hence $p_i(TM)=p_i(TW|_M)=j^\ast p_i(TW)$ for all $i$
(a trivial summand does not change Pontryagin classes, and Pontryagin classes are natural
under the restriction $j^\ast$).  Therefore, for any $\lambda\in P(k)$,
\[
p_\lambda[M]=\big\langle j^\ast\!\big(p_{\lambda_1}(TW)\smile\cdots\smile
p_{\lambda_r}(TW)\big),[M]\big\rangle .
\]
Writing $P:=p_{\lambda_1}(TW)\smile\cdots\smile p_{\lambda_r}(TW)\in H^{4k}(W)$ and using
$j_\ast[M]=0$ in $H_{4k}(W)$ together with the adjunction between $j^\ast$ and $j_\ast$
exactly as in Claim 1 of Theorem~\ref{thm:cobinv},
\[
p_\lambda[M]=\big\langle j^\ast P,[M]\big\rangle
=\big\langle P,\,j_\ast[M]\big\rangle=\big\langle P,0\big\rangle=0 .
\]
\end{proof}

\subsection*{From cobordism invariance to Fact (II)}

Theorem~\ref{thm:cobinv} and Proposition~\ref{prop:pnbdy} show that both functionals
$\tau$ and $p_\lambda$ ($\lambda\in P(k)$) are \emph{oriented-cobordism invariants}: they
vanish on every closed oriented $4k$--manifold that bounds, and they are additive under
disjoint union and (using also (I)) multiplicative under products.  Consequently they
descend to linear functionals on the rational oriented cobordism vector space
\[
V_k:=\Omega^{SO}_{4k}\otimes_{\mathbb Z}\mathbb Q ,
\]
the $\mathbb Q$--vector space whose elements are formal $\mathbb Q$--combinations of
closed oriented $4k$--manifolds modulo the relations ``$[M]=0$ if $M$ bounds'' and
``$[M\sqcup M']=[M]+[M']$''.  The Pontryagin numbers $\{p_\lambda\}_{\lambda\in P(k)}$ and
the signature $\tau$ are well-defined elements of the dual space $V_k^\ast$.

\begin{theorem}[Fact (II), conditional form]\label{thm:factII}
Suppose
\[
\dim_{\mathbb Q}V_k\;\le\;p(k).
\tag{$\sharp$}\label{eq:dimbound}
\]
Then Fact \textnormal{(II)} holds: there exist rationals $a_\lambda$
$(\lambda\in P(k))$ with $\tau(M)=\sum_\lambda a_\lambda p_\lambda[M]$ for every closed
oriented $4k$--manifold $M$.
\end{theorem}

\begin{proof}
Theorem~\ref{thm:span} produced $p(k)$ manifolds $M_\mu$ whose Pontryagin vectors
$\mathbf p[M_\mu]=(p_\lambda[M_\mu])_\lambda$ are linearly independent in
$\mathbb Q^{P(k)}$.  Since each $p_\lambda$ is a cobordism invariant, the linear
independence of these vectors forces the classes $[M_\mu]\in V_k$ to be linearly
independent (a linear relation among the $[M_\mu]$ would, on applying every $p_\lambda$,
give the same relation among the $\mathbf p[M_\mu]$).  Hence
$\dim_{\mathbb Q}V_k\ge p(k)$, and with \eqref{eq:dimbound} we get
$\dim_{\mathbb Q}V_k=p(k)$ and that $\{[M_\mu]\}_{\mu\in P(k)}$ is a basis of $V_k$.

Now the $p(k)$ functionals $\{p_\lambda\}_{\lambda\in P(k)}\subset V_k^\ast$ are linearly
independent: a relation $\sum_\lambda c_\lambda p_\lambda=0$ in $V_k^\ast$ evaluated on
the basis $\{[M_\mu]\}$ gives $\sum_\lambda c_\lambda p_\lambda[M_\mu]=0$ for all $\mu$,
i.e.\ the nonsingular matrix $A$ of Theorem~\ref{thm:span} annihilates $(c_\lambda)$,
whence all $c_\lambda=0$.  Thus $\{p_\lambda\}_{\lambda\in P(k)}$ is a basis of the
$p(k)$--dimensional space $V_k^\ast$.  In particular the cobordism invariant
$\tau\in V_k^\ast$ is a $\mathbb Q$--linear combination of the $p_\lambda$, say
$\tau=\sum_\lambda a_\lambda p_\lambda$ in $V_k^\ast$; evaluating on the class of an
arbitrary closed oriented $4k$--manifold $M$ gives
$\tau(M)=\sum_\lambda a_\lambda p_\lambda[M]$.  This is Fact (II).
\end{proof}

\begin{remark}[What remains, stated honestly]\label{rem:honest}
Theorem~\ref{thm:cobinv} and Proposition~\ref{prop:pnbdy} are proved here \emph{in full}
and use nothing beyond Poincar\'e--Lefschetz duality and linear algebra; they are the
substantive, self-contained step toward Fact (II) and already establish the lower bound
$\dim_{\mathbb Q}V_k\ge p(k)$.  The only ingredient \emph{not} reduced to elementary
manifold topology here is the matching upper bound \eqref{eq:dimbound},
$\dim_{\mathbb Q}V_k\le p(k)$.  This is the (rational) computation of the oriented
cobordism group, Thom's theorem
$\Omega^{SO}_\ast\otimes\mathbb Q=\mathbb Q[[\mathbb{CP}^2],[\mathbb{CP}^4],\dots]$.
By Theorem~\ref{thm:factII} the entire signature theorem follows from this single
classical input, and from no other deep machinery: no Atiyah--Singer index theorem, no
heat kernel, and no further use of cobordism beyond the rational rank
\eqref{eq:dimbound}.  We do not reprove \eqref{eq:dimbound} here; we have isolated it as
the precise and only remaining nonelementary hypothesis, and have shown that everything
else --- the cobordism invariance of $\tau$, the spanning family, the value
computations, and the deduction $\tau=L[\,\cdot\,]$ --- is elementary and complete.
\end{remark}

\bigskip

\section{Part (b): the functional equation of the $L$--genus and its geometric
meaning}\label{sec:partb}

\subsection*{The series, the genus, and the functional equation}

Hirzebruch's $L$--genus is the multiplicative genus attached to the even power series
\[
Q(z)=\frac{z}{\tanh z}=z\coth z
=1+\frac{z^2}{3}-\frac{z^4}{45}+\frac{2z^6}{945}-\cdots ,
\tag{10}\label{eq:Qseries}
\]
which is the function denoted $\sqrt{x}/\tanh\sqrt{x}$ in the problem statement under the
substitution $x=z^2$.  The ``virtual signature / $L$--genus'' associated with $Q$ is, for
a partition $\lambda$ describing formal roots, the value of the multiplicative sequence;
the basic analytic object behind it is the function $\coth$, which satisfies the
\emph{functional equation}
\[
\coth(z)+\coth(-z)=0,\qquad
\coth(z+w)=\frac{1+\coth z\,\coth w}{\coth z+\coth w},
\tag{11}\label{eq:funeq}
\]
together with the residue normalisation $\operatorname{Res}_{z=0}\coth z=1$, equivalently
$z\coth z\to 1$ as $z\to 0$, and the periodicity/pole structure
$\coth(z+\pi i)=\coth z$, with simple poles of residue $1$ at $z\in\pi i\,\mathbb Z$.  We
explain the geometric content of \eqref{eq:funeq} in three layers: (i) the residue/pole
normalisation, (ii) the addition law, and (iii) the resulting multiplicativity in fibre
bundles.

\subsection*{(i) Normalisation: the residue is the signature of a point of intersection}

The computation of $L[\mathbb{CP}^{2m}]$ in Proposition~\ref{prop:values} is, geometrically,
a residue at $z=0$:
\[
L[\mathbb{CP}^{2m}]=\big[h^{2m}\big]\,Q(h)^{2m+1}
=\frac{1}{2\pi i}\oint Q(z)^{2m+1}\,\frac{dz}{z^{2m+1}}
=\operatorname*{Res}_{z=0}\frac{(\coth z)^{2m+1}}{1}\,dz=1 .
\]
The substitution $w=\tanh z$ turns the pole of order $2m+1$ of $\coth^{2m+1}$ at $z=0$
into the pole $w^{-(2m+1)}$ and the volume form $dz$ into $dw/(1-w^2)$; the residue is the
coefficient of $w^{2m}$ in $(1-w^2)^{-1}=\sum_{j\ge0}w^{2j}$, namely $1$.  Geometrically,
this normalisation $\operatorname{Res}_{z=0}\coth z=1$ encodes that a single transverse
self-intersection point of a middle-dimensional cycle contributes $+1$ to the
intersection form on $\mathbb{CP}^{2m}$ (Proposition~\ref{prop:sig}); the genus is
calibrated so that complex projective spaces, the basic building blocks of positive
signature, all have $L$--genus $+1$.  Thus $\tau(\mathbb{CP}^{2m})=L[\mathbb{CP}^{2m}]=1$
is the geometric meaning of the residue normalisation.

\subsection*{(ii) The addition law: Novikov additivity and the cotangent two--variable
form}

The addition law in \eqref{eq:funeq} is the analytic shadow of the additivity of the
signature under gluing.  Recall \emph{Novikov additivity}: if a closed oriented
$4k$--manifold $M$ is cut along a separating closed hypersurface
$N^{4k-1}$ into $M=M_+\cup_N M_-$ (with $\partial M_\pm=N$, $M_+$ and $M_-$ glued along
$N$ with opposite induced orientations), then
\[
\tau(M)=\tau(M_+)+\tau(M_-),
\tag{12}\label{eq:novikov}
\]
where $\tau(M_\pm)$ is the signature of the intersection form on the image of
$H^{2k}(M_\pm,N)\to H^{2k}(M_\pm)$ (the form is well defined on this image by
Lefschetz duality, exactly as in Theorem~\ref{thm:cobinv}).  Equation
\eqref{eq:novikov} is proved by the same Poincar\'e--Lefschetz/Witt-index bookkeeping used
in Theorem~\ref{thm:cobinv}, applied to the Mayer--Vietoris decomposition of
$H^{2k}(M)$; the boundary term that would obstruct naive additivity is precisely a
totally isotropic piece, which contributes $0$ to the signature.  In the genus picture,
$\tau$ is computed as a sum of local residue contributions, one per ``building block'';
the additivity \eqref{eq:novikov} corresponds analytically to the fact that $\coth$ is an
\emph{odd} function with a simple pole of residue $1$ at the origin and is otherwise
regular, so residues of the relevant differential add when contributions are concatenated.
The two--variable form $\coth(z+w)=\dfrac{1+\coth z\coth w}{\coth z+\coth w}$ is the
algebraic record that $z\coth z$ is the unique normalised even genus for which these
local residue contributions combine consistently when blocks are stacked --- equivalently,
the unique genus whose value on a point is $1$ and which is \emph{strictly multiplicative}.

\subsection*{(iii) Multiplicativity in fibre bundles: the geometric face of the functional
equation}

The deepest geometric interpretation of the functional equation is the
\emph{multiplicativity of the signature in fibre bundles}, the Chern--Hirzebruch--Serre
phenomenon.  Among all multiplicative genera, the $L$--genus is singled out by the
following property, which is the geometric content of \eqref{eq:funeq}.

\begin{proposition}[Product multiplicativity, the elementary case]\label{prop:prodmult}
For closed oriented manifolds $F,B$ of dimensions divisible by $4$,
\[
\tau(F\times B)=\tau(F)\,\tau(B),\qquad L[F\times B]=L[F]\,L[B].
\]
\end{proposition}

\begin{proof}
This is fact (I), Equation~\eqref{eq:mult}: the intersection form of a product is the
tensor product of the factors' forms (K\"unneth), and the signature of a tensor product
of symmetric forms is the product of the signatures; while $L[\,\cdot\,]$ is multiplicative
because $Q$ defines a multiplicative sequence and $p(T(F\times B))=p(TF)\times p(TB)$.  The
agreement of the two multiplicative laws on the building blocks $\mathbb{CP}^{2m}$ (both
give $1$) is the residue normalisation of (i).
\end{proof}

\noindent
The functional equation \eqref{eq:funeq} expresses that the genus defined by $Q(z)=z\coth
z$ is the unique normalised genus that remains multiplicative not only for products but for
arbitrary oriented fibre bundles $F\to E\to B$ with structure group acting trivially on
$H^\ast(F;\mathbb Q)$, in which case
\[
\tau(E)=\tau(F)\,\tau(B).
\tag{13}\label{eq:CHS}
\]
The geometric mechanism is again a residue/localisation: integrating the $L$--form along
the fibre converts the fibrewise contribution into the constant $L[F]=\tau(F)$ times the
base form, precisely because $z\coth z$ is the generating function whose residues encode
fibre signatures additively under the monodromy action.  The two--variable cotangent
identity in \eqref{eq:funeq} is the infinitesimal/formal-group statement underlying
\eqref{eq:CHS}: $\coth$ is the exponential of the multiplicative formal group law
normalised so that the associated genus assigns the value $1$ to every point and is stable
under the fibration operation.  Equivalently, in Hirzebruch's residue formulation, the
signature is the sum over fixed data of local terms built from $\coth$, and the addition
formula \eqref{eq:funeq} is exactly the compatibility that makes this sum a cobordism
invariant equal to $\tau$.

\subsection*{Summary of the geometric interpretation}

The functional equation of the $L$--genus has the following layered geometric meaning.
\begin{itemize}
\item[(a)] The \emph{residue normalisation} $\operatorname{Res}_{z=0}\coth z=1$ /
$z\coth z|_{z=0}=1$ encodes that a transverse positive intersection point contributes
$+1$ to the signature; it is why every complex projective space has $L$--genus and
signature equal to $1$ (Proposition~\ref{prop:sig}, Proposition~\ref{prop:values}).
\item[(b)] The \emph{oddness} $\coth(-z)=-\coth z$ together with the simple-pole
structure is the analytic shadow of \emph{Novikov additivity} \eqref{eq:novikov}: signature
contributions of glued pieces add, and the boundary terms are totally isotropic, hence
signature-neutral (this is the same Witt-index mechanism as in
Theorem~\ref{thm:cobinv}).
\item[(c)] The \emph{addition law} for $\coth$ is the formal-group / two--variable
identity underlying the \emph{multiplicativity of the signature in fibre bundles}
\eqref{eq:CHS}, of which the product formula
$\tau(F\times B)=\tau(F)\tau(B)$ (Proposition~\ref{prop:prodmult}) is the elementary
special case.  Among all normalised genera, $z\coth z$ is the unique one whose residues
combine consistently under these geometric operations, which is why it --- and no other
--- computes the signature.
\end{itemize}
This completes the geometric interpretation requested in part (b).


\end{document}
